\section{Other Refinements, Intermediate Results and Limits}\label{lim:section}

This section is not needed for Theorem~\ref{thm:main}. It records the
refinements behind the intermediate exponents of
Table~\ref{intro:table-steps}, with proofs, and what is known about the limits
of the method. Statements about limits are labeled by the kind of evidence
behind them. We use the notation of Section~\ref{an:section}: $\mathcal K$,
$F$, $d$, $\theta$, the residue degree $f_K(\mathfrak p)$ of a prime
$\mathfrak p$ of $B$ in $K/B$, the free and selected primes
(Definition~\ref{an:selected-def}), $g_s$ and $w$
(Subsection~\ref{an:budget-section}), and $\kappa_\infty$, $Z_{\rm sel}$ and
$T_{\rm sel}$.

\subsection{The adaptive dichotomy}
The comparison with a fixed field charges each free prime $\mathfrak p$ of $B$
with its floor $f^0_{\mathfrak p}$, a lower bound for its residue degree
$f_K(\mathfrak p)$. Where $f_K(\mathfrak p)$ is not known, an adversary
choosing the fields cannot both keep it small and avoid a gain: if
$f_K(\mathfrak p)$ is small, the places of $F$ above $\mathfrak p$ that split
in $K$ can carry shell profiles; if it is large, the term of $\mathfrak p$ in
$\zeta_K$ is smaller; and places of $F$ above $\mathfrak p$ that are inert in
$K$ lower $L_F$. This subsection makes the three gains precise
(Definition~\ref{lim:credits}) and proves that their minimum can be added to
the margin (Propositions~\ref{lim:A} and~\ref{lim:B}). With the bound
$C=\ceiling$ alone, Proposition~\ref{lim:B} gave the exponent $1.042925$.

\begin{lemma}\label{lim:inert}
Let $K\in\mathcal K$ and $s>1$ real. Then
\[
 \frac1d\log L_F(s)=\frac1{[K:\Q]}\log\zeta_K(s)-\frac1d\sum_{\mathfrak q}
 \artanh(\mathrm N\mathfrak q^{-s}),
\]
the sum over the primes $\mathfrak q$ of $F$ that are inert in $K$. For a free
prime $\mathfrak p$ of $B$ with $f=f_K(\mathfrak p)$ and $N=\mathrm N\mathfrak
p$, let $\rho_K(\mathfrak p)$ be the proportion of the primes $\mathfrak P$ of
$K$ above $\mathfrak p$ with $\iota_1\mathfrak P=\mathfrak P$. The primes of
$F$ above $\mathfrak p$ that are inert in $K$ are the $\mathfrak P\cap F$ with
$\iota_1\mathfrak P=\mathfrak P$; there are $\rho_K(\mathfrak p)d/f$ of them,
each of norm $N^{f/2}$, and $\rho_K(\mathfrak p)>0$ only if $f$ is even. The
primes of $F$ above $\mathfrak p$ that split in $K$ number
$(1-\rho_K(\mathfrak p))d/(2f)$, each of norm $N^f$.
\end{lemma}

\begin{proof}
At a prime $\mathfrak q$ of $F$, with $x=\mathrm N\mathfrak q^{-s}$, the Euler
factor of $L_F$ is $(1-x)^{-1}$ if $\mathfrak q$ splits in $K$ and
$(1+x)^{-1}$ if $\mathfrak q$ is inert; no finite prime of $F$ ramifies in $K$
(G2). The factor of $\zeta_K$ is $(1-x)^{-2}$, respectively
$(1-x^2)^{-1}$, and $-\frac12\log(1-x^2)+\log(1+x)=\artanh x$. Since
$[K:\Q]=2d$ this gives the identity. The prime $\mathfrak p$ is unramified in
$K$, and $K/B$ is Galois of degree $d$, so there are $d/f$ primes $\mathfrak
P$ above $\mathfrak p$, with cyclic decomposition groups of order $f$. The
primes of $K$ above $\mathfrak q=\mathfrak P\cap F$ are $\mathfrak P$ and
$\iota_1\mathfrak P$. If $\iota_1\mathfrak P=\mathfrak P$, then $\mathfrak q$
is inert, $\mathrm N\mathfrak q=N^{f/2}$, and $\iota_1$ lies in the cyclic
decomposition group of $\mathfrak P$, so $f$ is even. If
$\iota_1\mathfrak P\ne\mathfrak P$, then $\mathfrak q$ splits and
$\mathrm N\mathfrak q=N^f$.
\end{proof}

\begin{definition}\label{lim:credits}
Fix $\sigma=1+\varepsilon\in(1,2]$, $0<\delta<1$, and floors $f^0_{\mathfrak
p}$ for the free primes, that is, integers with $f^0_{\mathfrak p}\mid
f_K(\mathfrak p)$ for every $K\in\mathcal K$. For a free prime $\mathfrak p$
with $N=\mathrm N\mathfrak p$, and for $f\in f^0_{\mathfrak p}\cdot2^{\mathbb
N}$, put
\[
 \begin{gathered}
 R_{\mathfrak p}(f)=\frac{g_\sigma(N^{f^0_{\mathfrak p}})}{2f^0_{\mathfrak p}}
 -\frac{g_\sigma(N^f)}{2f}\ge0,
 \qquad
 I_{\mathfrak p}(f)=\begin{cases}\artanh(N^{-f/2})/f&f\text{ even},\\0&f\text{ odd},\end{cases}\\
 S_{\mathfrak p}(f)=\frac1{2f}\max\Bigl(0,\max_{k\ge1}\bigl(\log(k+1)-\delta kf\log N\bigr)\Bigr),
 \end{gathered}
\]
with $R_{\mathfrak p}(\infty)=g_\sigma(N^{f^0_{\mathfrak p}})/(2f^0_{\mathfrak
p})$ and $I_{\mathfrak p}(\infty)=S_{\mathfrak p}(\infty)=0$, and let
\[
 G_{\mathfrak p}=\inf\bigl\{R_{\mathfrak p}(f)+\rho I_{\mathfrak p}(f)
 +(1-\rho)S_{\mathfrak p}(f)\bigr\}\ge0,
\]
over $f\in f^0_{\mathfrak p}\cdot2^{\mathbb N}\cup\{\infty\}$ and
$\rho\in[0,1]$, with $\rho=0$ if $f$ is odd. For $K\in\mathcal K$ we call
$R_{\mathfrak p}(f_K(\mathfrak p))+\rho_K(\mathfrak p)I_{\mathfrak
p}(f_K(\mathfrak p))$ the \emph{adaptive credit} of $\mathfrak p$ (it lowers
the bound for $d^{-1}\log L_F(1)$, Proposition~\ref{lim:A}), and
$(1-\rho_K(\mathfrak p))S_{\mathfrak p}(f_K(\mathfrak p))$ its \emph{window
credit} (it raises the finite term of the margin, Proposition~\ref{lim:B}).
\end{definition}

\begin{lemma}\label{lim:credit-sum}
In Definition~\ref{lim:credits}, $R_{\mathfrak p}(f)\ge0$. For every
$K\in\mathcal K$ and every free prime $\mathfrak p$, the sum of the adaptive
credit and the window credit of $\mathfrak p$ is at least $G_{\mathfrak p}$.
\end{lemma}

\begin{proof}
The inequality $R_{\mathfrak p}(f)\ge0$ follows from the monotonicity in
Lemma~\ref{an:local-contribution}. Since $f_K(\mathfrak p)$ is a power of two
divisible by $f^0_{\mathfrak p}$, and $\rho_K(\mathfrak p)=0$ if
$f_K(\mathfrak p)$ is odd (Lemma~\ref{lim:inert}), the sum of the adaptive
credit and the window credit of $\mathfrak p$ is one of the expressions in the
infimum defining $G_{\mathfrak p}$, and hence at least $G_{\mathfrak p}$.
\end{proof}

\begin{proposition}[Adaptive ceiling]\label{lim:A}
Fix $\sigma$, $\delta$ and floors as in Definition~\ref{lim:credits}, with
the functions $R_{\mathfrak p}$ and $I_{\mathfrak p}$ defined there. Let
$Y_*$ be a real number such that, for every $K\in\mathcal K$ and every
finite set $\mathcal A$ of free primes,
\begin{equation}\label{lim:termwise}
 \frac{\log\zeta_K(\sigma)}{[K:\Q]}\le Y_*-\sum_{\mathfrak p\in\mathcal A}
 R_{\mathfrak p}(f_K(\mathfrak p)),
\end{equation}
let $Y'_*$ be as in Proposition~\ref{an:prime-budget}, and let
$C'>Y_*+\varepsilon(\kappa_\infty+Y'_*)$. For every finite set $\mathcal A$
of free primes there is an $m_0$ such that every $K\in\mathcal K$ with
$[K:B]\ge m_0$ satisfies
\[
 \frac1d\log L_F(1)<C'-\sum_{\mathfrak p\in\mathcal A}\bigl(R_{\mathfrak p}(f_K(\mathfrak p))
 +\rho_K(\mathfrak p)I_{\mathfrak p}(f_K(\mathfrak p))\bigr),
\]
that is, $C'$ minus the adaptive credits of the primes of $\mathcal A$.
\end{proposition}

\begin{remark}\label{lim:termwise-remark}
The hypothesis \eqref{lim:termwise} holds for the termwise bounds of this
paper: by \eqref{an:local-sum}, $Y_*$ is a sum of terms
$g_\sigma(\mathrm N\mathfrak p^{f^0_{\mathfrak p}})/(2f^0_{\mathfrak p})$ over the
free primes, plus terms for the selected primes, and keeping the exact term
$g_\sigma(N^f)/(2f)$ at the primes of $\mathcal A$ lowers it by
$\sum_{\mathcal A}R_{\mathfrak p}$.
\end{remark}

\begin{proof}[Proof of Proposition~\ref{lim:A}]
Suppose not, and let $K^{(j)}$ be fields of strictly increasing degree that
violate the inequality. As $\mathcal A$ is finite and the
$f_j(\mathfrak p)=f_{K^{(j)}}(\mathfrak p)$ are powers of two times
$f^0_{\mathfrak p}$, we may pass to a subsequence along which, for each
$\mathfrak p\in\mathcal A$, either $f_j(\mathfrak p)=f_\infty(\mathfrak p)$ is
constant or $f_j(\mathfrak p)\to\infty=:f_\infty(\mathfrak p)$, and
$\rho_j(\mathfrak p)=\rho_{K^{(j)}}(\mathfrak p)\to\rho_\infty(\mathfrak p)$;
and then, as in Lemma~\ref{an:limit}, to one along which the $\beta_q$ exist,
so that (a) and (b) of that lemma hold. Then $R_{\mathfrak p}(f_j)\to
R_{\mathfrak p}(f_\infty)$ and $\rho_jI_{\mathfrak p}(f_j)\to\rho_\infty
I_{\mathfrak p}(f_\infty)$.

Fix $s\in(1,2]$. By the proof of Lemma~\ref{an:shift},
$d^{-1}\log L_F(1)\le d^{-1}\log L_F(s)+\Upsilon(s-1,\theta)$, and by
Lemma~\ref{lim:inert}, dropping the inert primes above primes outside
$\mathcal A$,
\[
 \frac1d\log L_F(s)\le\frac{\log\zeta_K(s)}{[K:\Q]}-\sum_{\mathfrak p\in\mathcal A}
 \frac{\rho_K(\mathfrak p)}{f_K(\mathfrak p)}
 \artanh\bigl(\mathrm N\mathfrak p^{-f_K(\mathfrak p)s/2}\bigr).
\]
Along the subsequence the right side tends to $Z(s)-I_\infty(s)$, with $Z$ as
in Lemma~\ref{an:limit} and
\[
 I_\infty(s)=\sum_{\mathfrak p\in\mathcal A}\frac{\rho_\infty}{f_\infty}
 \artanh\bigl(\mathrm N\mathfrak p^{-f_\infty s/2}\bigr),
\]
a term being zero when $f_\infty=\infty$. Hence
\[
 C'-\sum_{\mathfrak p\in\mathcal A}\bigl(R_{\mathfrak p}(f_\infty)+\rho_\infty
 I_{\mathfrak p}(f_\infty)\bigr)\le Z(s)-I_\infty(s)+\sup_\theta\Upsilon(s-1,\theta).
\]
As $s\downarrow1$, $Z(s)\to Z(1)$, $I_\infty(s)\to\sum_{\mathfrak p\in\mathcal A}
\rho_\infty I_{\mathfrak p}(f_\infty)$ and the supremum tends to $0$; so
$C'-\sum_{\mathcal A}R_{\mathfrak p}(f_\infty)\le Z(1)$. As in the proof of
Proposition~\ref{an:prime-budget},
$Z(1)\le Z(1+\varepsilon)+\varepsilon(\kappa_\infty+Y'_*)$, and by
\eqref{lim:termwise} and Lemma~\ref{an:limit}(a),
$Z(1+\varepsilon)\le Y_*-\sum_{\mathcal A}R_{\mathfrak p}(f_\infty)$. Together,
$C'-\sum_{\mathcal A}R_{\mathfrak p}\le Y_*+\varepsilon(\kappa_\infty+Y'_*)
-\sum_{\mathcal A}R_{\mathfrak p}<C'-\sum_{\mathcal A}R_{\mathfrak p}$, a
contradiction.
\end{proof}

\begin{lemma}\label{lim:transfer-C}
Proposition~\ref{fw:transfer} remains true if its hypothesis
$\log L_{F_j}(1)\le d_jC$ is replaced by $\log L_{F_j}(1)\le d_jC_j$, and $C$ by
$C_j$ in the margin \eqref{fw:margin}.
\end{lemma}

\begin{proof}
In the proof of Proposition~\ref{fw:transfer} the hypothesis enters only once,
for the field at hand, in the lower bound for the first factor through
\eqref{geo:mass-density}. The choice of $\eta$ and the Chebyshev step depend
only on the infimum of the margins and on the finite set of parameters.
\end{proof}

\begin{proposition}[Adaptive margin]\label{lim:B}
Let $\mathcal A$ be a finite set of free primes, let $C'$ satisfy the
hypotheses of Proposition~\ref{lim:A}, and let $G_{\mathfrak p}$ be as in
Definition~\ref{lim:credits}. Let $\mathcal P$ and $\mathcal R$ be as
in Definition~\ref{cert:data}(e), with the effective types $(e_r,f_r)$ of
Definition~\ref{cert:effective-types}, and for $r\in\mathcal R$ let $g_r$ be a
shell profile with parameters $(Q_r,k_r,w^{(r)})$, placed at the places of
$\mathcal P$ above $r$, where $Q_r$ and $k_r$ are those of
Table~\ref{cert:tab-local} and only the weights $w^{(r)}$ may differ from those
of Definition~\ref{cert:data}. Let $(f_\R,g_\C)$ be an admissible pair whose Fourier
constants satisfy \eqref{fw:fourier-condition} with $H_f=H$, the number $H$ of
Lemma~\ref{cert:places}. Let $J_\C^*\le J_\C$, and let $\mathcal M_*(\theta)$
be defined by the formula of Definition~\ref{cert:lower-bounds} from these
data, with $C'$ in place of $C$. If the slope $-\log\pi-2J_\R+J_\C^*$ of
$\mathcal M_*$ in $\theta$ is positive and
$\mathcal M_*(\theta_*)+\sum_{\mathfrak p\in\mathcal A}G_{\mathfrak p}>0$, then
there are finite sets $U_j\subset\R^2$ with $|U_j|\to\infty$ and
$u(U_j)/|U_j|^{1+\delta}\to\infty$.
\end{proposition}

\begin{proof}
Choose $K_j\in\mathcal K$ with $[K_j:B]\to\infty$ and $[K_j:B]\ge m_0$ of
Proposition~\ref{lim:A}. For $\mathfrak p\in\mathcal A$ write
$f_{K_j}=f_{K_j}(\mathfrak p)$ for the residue degree of $\mathfrak p$ in
$K_j/B$ (not to be confused with the residue degrees $f_r$ of the effective
types or with the profile $f_\R$) and $\rho_j=\rho_{K_j}(\mathfrak p)$, and
put $C_j=C'-\sum_{\mathfrak p\in\mathcal
A}(R_{\mathfrak p}(f_{K_j})+\rho_jI_{\mathfrak p}(f_{K_j}))$, so that $\log
L_{F_j}(1)\le d_jC_j$. Let $\mathcal P_j$ consist of the places of
$\mathcal P$ for $K_j$ and, for each $\mathfrak p\in\mathcal A$ with
$S_{\mathfrak p}(f_{K_j})>0$, of the places of $F_j$ above $\mathfrak p$ that
split in $K_j$, with the unweighted shell profile with parameters
$(\mathrm N\mathfrak p^{f_{K_j}},k)$, where $k$ attains the maximum in
$S_{\mathfrak p}(f_{K_j})$. By Lemma~\ref{lim:inert} these places add the window
credit $(1-\rho_j)S_{\mathfrak p}(f_{K_j})$ to
$d_j^{-1}\sum\log\mathcal F_{\delta,\mathfrak p}$. Since
$S_{\mathfrak p}(f)=0$ for large $f$, their parameters lie in a finite set. The
Fourier condition still holds, because $H_f$ only increases. By
Lemma~\ref{cert:places} and $J_\C\ge J_\C^*$, the margin of $K_j/F_j$ is at
least $\mathcal M_*(\theta_j)+\sum_{\mathcal A}(R_{\mathfrak p}+\rho_j
I_{\mathfrak p}+(1-\rho_j)S_{\mathfrak p})\ge\mathcal M_*(\theta_*)+\sum_{\mathcal A}
G_{\mathfrak p}>0$, by the positive slope and Lemma~\ref{lim:credit-sum}, and
Lemma~\ref{lim:transfer-C} applies.
\end{proof}

\subsection{A refined Tsfasman--Vl\u{a}du\c{t} step with one abscissa}

This subsection records the one-abscissa version of
Proposition~\ref{an:signed}, combined with the adaptive credits, used for the
intermediate exponent $1.043116$.

\begin{proposition}\label{lim:D}
Let $a\ge1$ and $b_{\rm TV}\ge0$ satisfy
\begin{equation}\label{lim:kernel}
 g_1(q)\le a\,g_\sigma(q)+b_{\rm TV}\,w(q)\qquad\text{for every real }q\ge2.
\end{equation}
Then Proposition~\ref{lim:A}, and hence Proposition~\ref{lim:B}, hold with any
$C'>Z_{\rm sel}(1)+a\bigl(Y_*-Z_{\rm sel}(\sigma)\bigr)+b_{\rm TV}\,(2\kappa_\infty-T_{\rm sel})$
in place of the condition $C'>Y_*+\varepsilon(\kappa_\infty+Y'_*)$.
\end{proposition}

\begin{proof}
In the proof of Proposition~\ref{lim:A}, keep everything up to
$C'-\sum_{\mathcal A}R_{\mathfrak p}(f_\infty)\le Z(1)$. Write $\beta^{\rm
free}=\beta-\beta^{\rm sel}$ as in the proof of Proposition~\ref{an:signed}.
By \eqref{lim:kernel}, Lemma~\ref{an:limit}(b) and $b_{\rm TV}\ge0$,
\[
 Z(1)=Z_{\rm sel}(1)+\sum_q\beta^{\rm free}_qg_1(q)
 \le Z_{\rm sel}(1)+a\bigl(Z(\sigma)-Z_{\rm sel}(\sigma)\bigr)+b_{\rm TV}\,(2\kappa_\infty-T_{\rm sel}).
\]
By \eqref{lim:termwise}, $Z(\sigma)\le Y_*-\sum_{\mathcal A}R_{\mathfrak
p}(f_\infty)$, and since $a\ge1$ and $R\ge0$ (Lemma~\ref{lim:credit-sum}), $a\sum_{\mathcal A}R\ge
\sum_{\mathcal A}R$. So $C'-\sum_{\mathcal A}R\le Z_{\rm sel}(1)+a(Y_*-Z_{\rm
sel}(\sigma))+b_{\rm TV}(2\kappa_\infty-T_{\rm sel})-\sum_{\mathcal A}R$, a
contradiction. The bound $Y'_*$ is not needed.
\end{proof}

\begin{remark}\label{lim:kernel-check}
For $q\ge q_{\min}=41^4$, condition \eqref{lim:kernel} follows from $h(x)\ge0$
for $x=\log q$, where $h(x)=ae^{-\varepsilon x}+2b_{\rm TV}x(1-e^{-x})-1-2e^{-x}$,
by $g_1(q)\le1/(q-1)$, $g_\sigma(q)\ge q^{-\sigma}$, $w(q)\ge2\log q/(q+1)$,
$q/(q+1)\ge1-1/q$ and $q/(q-1)\le1+2/q$; the program \texttt{dihedral/lptv.py}
checks $h\ge0$ in Arb on cells, with bisection near the minimum, and in a
closed form, and checks \eqref{lim:kernel} directly on $5950$ geometric cells
of $[2,2q_{\min}]$. Proposition~\ref{an:signed} with $n=1$ and $c_1=a$ is this
step without the adaptive credits; with several abscissae and coefficients of
both signs it gains $1.18\cdot10^{-4}$ in $C$ on the same data.
\end{remark}

\subsection{Intermediate results}
This subsection lists how the intermediate exponents of
Table~\ref{intro:table-steps} were obtained, as special cases of the results
above. Each has a witness file and a replay program in the directory
\texttt{certificates} of \cite{NaslundCode4317}.
\begin{enumerate}
\item[(1)] $1.042925$: Proposition~\ref{lim:B} with $C'=\ceiling$, the
 floors relative to $E$ of Remark~\ref{an:base-ceiling}, $\sigma=301/300$, and
 $\mathcal A$ the free primes of norm at most $X=1.2\cdot10^7$
 (\texttt{reproduce\_adaptive.py}; the values of $\sigma$ and $X$ are in the
 witness file \texttt{witness\_0.042925.json}).
\item[(2)] $1.042965$: the census of Proposition~\ref{dh:census} up to
 $10^{12}$ relative to $E$, which raises the floors of the detected primes of
 vector $0$ from $1$ to $2$ (\texttt{reproduce\_census.py}).
\item[(3)] $1.043113$ to $1.043124$: Lemma~\ref{dh:EW} for the subspace
 $W=\langle\psi_{10},\psi_{23}\rangle\subset W''$ in place of $W''$
 (Remark~\ref{dh:EW-sub}), whose nonzero elements
 $\psi_{10},\psi_{23},\psi_{24}$ are all of dihedral type, so that
 $[E_W:\Q]=2048$ and $\zeta_{E_W}$ is $\zeta_E$ times the squares of the $192$
 $L$-functions of these three forms; with the census relative to $E_W$ up to
 $8.5\cdot10^{14}$, and with Propositions~\ref{an:prime-budget}, \ref{lim:D}
 and~\ref{an:signed} in turn (\texttt{dihedral/reproduce\_dihedral.py}). One
 review found that an earlier version of this step had credited the thirteen
 inert primes of vector $0$ (Lemma~\ref{dh:inert-v0}) twice; the corrected value is in the table.
\item[(4)] $1.043152$: Lemma~\ref{dh:EW} for the subspace
 $W'=\langle\psi_{10},\psi_{23},\psi_{19}\rangle\subset W''$
 (Remark~\ref{dh:EW-sub}), whose seven nonzero elements
 are five forms of dihedral type and $\psi_{19}+\psi_{23}$,
 $\psi_{19}+\psi_{24}$ of rank $4$; $[E_{W'}:\Q]=4096$, with $32$
 $L$-functions of degree $8$ (\texttt{dihedral/reproduce\_w3.py}).
\end{enumerate}

\subsection{Limits}
This subsection records what is known about how far the method can go, with
the kind of evidence for each statement. The margin decreases by about $23.8$ per unit of $\delta$ and by $1$ per unit
of $C$, so $2.38\cdot10^{-5}$ in $C$ corresponds to $10^{-6}$ in $\delta$.

\emph{The next step (computation).} With general shell weights and
profiles optimized at $\delta=0.043172$, the margin is positive only if
$C\le0.0422631$. The repository contains, unused by any witness, the census
relative to $\EW$ extended to the limit $2^{50}\approx1.13\cdot10^{15}$ of the
program \texttt{census\_kv4.c}, and $35$ further abscissae. With both the linear programme of
Section~\ref{ce:section} reaches only about $0.0422638$, about $7\cdot10^{-7}$
short.

\emph{More dihedral forms (proved).} By Lemma~\ref{dh:W-structure}, every
space properly containing $W''$ contains a form whose polar form has rank at
least $6$, whose irreducible representations have dimension at least $8$ and
whose $L$-functions have degree at least $16$ over $\Q$, with conductors
around $10^{40}$. So $W''$ is the end of this route. The remaining part of
$C_{\rm eff}-Z_{\rm sel}(1)\approx6.1\cdot10^{-4}$ comes from the primes of
vector $0$ in $\ker W''$ that the census does not reach; the census gains only
like the logarithm of its range.

\emph{This design (proved, given the margin formula; numerical for the
profiles).} By Remark~\ref{an:limit-remark}, no bound of the kind of
Proposition~\ref{an:signed} is below $Z_{\rm sel}(1)=0.0416685$. With
$C=Z_{\rm sel}(1)$ and profiles optimized by our programs, the margin is about
$+3.9\cdot10^{-4}$ at $\delta=0.04318$ and $-8.8\cdot10^{-5}$ at
$\delta=0.04320$. So for this tower, these local types and this margin formula,
the method cannot go beyond about $\delta=0.043196$; this assumes that the
profile optimization is close to optimal.

\emph{A fourth cap (evidence).} One more $C_4$ cap would be worth about
$2.9\cdot10^{-4}$ in $\delta$, but it raises the value of the
Golod--Shafarevich function at its minimum from $-0.00211$ to $+0.00288$, so
Proposition~\ref{tw:infinite} no longer applies. Screens of weighted versions
of the Golod--Shafarevich argument, of the index-two subgroups of $G_B$, and
of further overlaps between the local modules did not recover the deficit.
These are limits of the present arguments, not proofs that such a tower is
finite.

\emph{A larger conjugacy class (conditional).} The proof uses
$\theta\ge\theta_*=1/2-2^{-17}$, from the class of $\iota_1$ in
$\overline G_B$, of size $2^{15}$ (Lemma~\ref{tw:retained-quotient}(c)). A class of
size $2^{16}$ in a suitable quotient would give $\theta\ge1/2-2^{-18}$ and add
about $2.4\cdot10^{-6}$ to the margin; with the bound $C\le0.04226506$, which
\texttt{signed3.py} certifies for $\EW$ with the census to $8.5\cdot10^{14}$ and
$36$ abscissae (not packaged as a witness), this would give the exponent
$1.043172$. An argument for such a class, by an agent run of the Muse Spark
system, is recorded in the directory \texttt{research/2026-10-muse-involution-class}
of \cite{NaslundCode4317}. It has been checked only inside that run and is not
used here.
